\documentclass[tikz]{standalone}
\usetikzlibrary{calc}

\begin{document}

\usepgflibrary{shadings} 

\pgfdeclareradialshading{myflameshade}{\pgfpoint{0}{0.05cm}}{%
  color(0.0cm)=(white); color(0.5cm)=(red!40!yellow);
  color(0.7cm)=(red!50!yellow); color(0.9cm)=(white)
}

\tikzset{
  vial/.cd,
  water height/.initial=0.74,
  water color/.initial=cyan!60!blue!20!white,
  flame scale/.initial=0.15,
  locked/.is choice,
  locked/true/.code={\def\islocked{1}},
  locked/false/.code={\def\islocked{0}},
  locked=false,
  %
  /tikz/flame/.pic={
    \def\wflame{1.0}\def\hflame{6.0}
    % transform shape here ensures the flame itself can be scaled
    \begin{scope}[transform shape]
      \clip (0,-0.1*\hflame) to[out=0, in=-90] (\wflame,0) 
            arc (0:180:\wflame\space and \hflame) to[out=-90,in=180] cycle;
      \shade[shading=myflameshade] (0,0) ellipse ({\wflame} and {\hflame});
    \end{scope}
  },
  %
  /tikz/vial/.pic={
    \tikzset{vial/.cd, #1}
    \pgfkeysgetvalue{/tikz/vial/water height}{\ratio}
    \pgfkeysgetvalue{/tikz/vial/water color}{\vcol}
    \pgfkeysgetvalue{/tikz/vial/flame scale}{\fscale}
    
    % Geometry Constants
    \def\Rx{1.1} \def\Ry{0.2} \def\H{2.6}
    \def\h{\ratio*\H}
    \def\pt{0.05*\H} \def\prx{0.040*\H} \def\pry{0.006*\H} \def\pl{0.25*\H}
    
    % transform shape is the magic key here
    \begin{scope}[transform shape]
        \colorlet{mydarkblue}{blue!50!black}
        \colorlet{darkwatercol}{\vcol}
        \colorlet{metalcol}{blue!40!black!10!white}
        \tikzstyle{dark water}=[draw=mydarkblue, top color=darkwatercol, bottom color=darkwatercol!80!black]
        \tikzstyle{metal}=[draw=metalcol!20!black, top color=metalcol, bottom color=metalcol!90!black]

        % The nested flame now inherits the vial's scale
        \pic[scale=\fscale] at (0, -1) {flame};
        
        \draw[dark water] (-\Rx,\h) -- (-\Rx,0) arc (180:360:{\Rx} and {\Ry}) -- (\Rx,\h);
        \draw[metal] (-\Rx,\h) --++ (0,-\pt) arc (180:360:{\Rx} and {\Ry}) --++ (0,\pt);
        \draw[metal] (0,\h) ellipse ({\Rx} and {\Ry});
        \draw[metal] (-\prx,\h+\pl) --++ (0,-\pl) arc (180:360:{\prx} and {\pry}) --++ (0,\pl);
        \draw[metal](0,\h+\pl) ellipse ({\prx} and {\pry});
        \draw[thick] (-\Rx,\H) -- (-\Rx,0) arc (180:360:{\Rx} and {\Ry}) -- (\Rx,\H) (0,\H) ellipse ({\Rx} and {\Ry});


        \ifnum\islocked=1
          \fill[black!70] (-\Rx-0.1, \h+0.1) rectangle (-\Rx+0.1, \h+0.2);
          \fill[black!70] (\Rx-0.1, \h+0.1) rectangle (\Rx+0.1, \h+0.2);
        \fi


    \end{scope}
  }
}

\begin{tikzpicture}

  \pic[scale=1.3] at (0,0) {vial={water height=0.4}};
  \node at (0,0.4) {$\rho(t_0),T(t_0),p$};

  \pic[scale=1.3] at (4,0) {vial={water height=0.6, water color=orange!20!white}};
  \node at (4,0.8) {$\rho(t_1),T(t_1),p$};

  \pic[scale=1.3] at (8,0) {vial={water height=0.8, water color=red!80!orange!40!white}};
  \node at (8,1.2) {$\rho(t_2),T(t_2),p$}; 
   
\end{tikzpicture}
\end{document}
